
\subsection{微分分次 c.f.z.}
\label{subsec:6-3-2}

现在考虑分次 $\Lambda_{0,\mathrm{nov}}$-模的情形。

\begin{definition}{6.3.8}\label{def:6-3-8}
设 $\widehat C$ 是分次
$\Lambda_{0,\mathrm{nov}}$-模，
并假设对每个 $k$，
$\widehat C^k$ 都是 c.f.z.。
\emph{微分分次 c.f.z.（d.g.c.f.z.）}\liftIndex{I334}
是一个偶对 $(\widehat C,\delta)$，
其中 $\delta:\widehat C\to\widehat C$
是次数为 $1$ 的算子，满足
\[
\delta\circ\delta=0,
\qquad
\delta(F^\lambda\widehat C)
\subseteq F^\lambda\widehat C.
\]
若每个 $\widehat C^k$
都是有限 c.f.z.，
则称这个偶对为\emph{有限 d.g.c.f.z.}。
\end{definition}

注意，$\widehat C^k$ 是
$\Lambda_{0,\mathrm{nov}}^{(0)}$-模。

下面证明一个对本节所构造谱序列的
某些收敛性质至关重要的命题。

\begin{proposition}{6.3.9}\label{prop:6-3-9}
设 $W$ 是 $\widehat C^k$ 的有限生成
$\Lambda_{0,\mathrm{nov}}^{(0)}$-子模。
则存在仅依赖于 $W$、与 $\lambda$ 无关的常数 $c$，使
\[
\delta(W)\cap F^\lambda\widehat C^{k+1}
\subseteq
\delta\bigl(W\cap F^{\lambda-c}\widehat C^k\bigr).
\]
\end{proposition}

\begin{proof}
显然，$\delta(W)$ 是 $\widehat C^{k+1}$
的有限生成子模，
故由引理~\ref{lem:6-3-2} 有标准基。
设 $\{\mathbf v_i\}$ 是它的一组标准基，
并取 $\mathbf w_i\in W$，使
$\delta\mathbf w_i=\mathbf v_i$。
设非零元 $\mathbf u\in\delta(W)$
满足 $E(\mathbf u)=\lambda$，并写作
\[
\mathbf u=\sum_i a_i\mathbf v_i.
\]
令 $\sigma(a_i)\in R$ 表示
$a_i\in\Lambda_{0,\mathrm{nov}}^{(0)}$
的首系数，
并记 $E(\mathbf v_i)=\lambda_i$，
$E(a_i)=\mu_i\ge0$。
注意，对每个 $a_i\ne0$，
都有 $\lambda\ge\lambda_i$。

\pagebreak
\begin{lemma}{6.3.10}\label{lem:6-3-10}
若对所有出现的 $i$ 有 $\lambda>\lambda_i$，
则 $\lambda_i+\mu_i\ge\lambda$。
\end{lemma}

\begin{proof}
反设某个 $i$ 满足
$\lambda_i+\mu_i<\lambda$。
令
\[
\lambda-a
=\inf_i\{\lambda_i+\mu_i\},
\qquad a>0.
\]
因为 $E(\mathbf u)=\lambda>\lambda-a$，
便有
\[
\sum_{i:\,\lambda-a=\lambda_i+\mu_i}
\sigma(a_i)\sigma(\mathbf v_i)=0.
\]
这与 $\mathbf v_i$ 的首系数
线性无关相矛盾。
引理得证。
\end{proof}

由引理~\ref{lem:6-3-10}，
若 $\lambda>\lambda_i$ 对所有
$a_i\ne0$ 成立，
则 $\mu_i=E(a_i)\ge\lambda-\lambda_i$。
现在取 $c=\max_i\{\lambda_i\}$；
它只依赖于 $W$，而与 $\lambda$ 无关。
若 $\lambda>c$，
则有 $b_i\in\Lambda_{0,\mathrm{nov}}^{(0)}$，
使 $T^{\lambda-c}b_i=a_i$。
因此
\[
\mathbf u
=\delta\Bigl(
\sum_i T^{\lambda-c}b_i\mathbf w_i
\Bigr)
\in
\delta\bigl(W\cap F^{\lambda-c}\widehat C^k\bigr).
\]
另一方面，若 $\lambda\le c$，
则 $F^{\lambda-c}\widehat C^k=\widehat C^k$，
命题显然成立。\translatornote{原书实体页 379
在最后一种情形写成
“若某个 $i$ 满足 $\lambda=\lambda_i$，
则 $\lambda=c$”。
由 $c=\max_i\lambda_i$ 只能推出
$\lambda\le c$；
这已足以证明此处的包含关系，
见疑误 A-0364。}
\end{proof}

\begin{lemma}{6.3.11}\label{lem:6-3-11}
设 $V$ 是有限 c.f.z.，
$W$ 是它的子模，且满足：
对任意 $\mathbf w\in W$，
若 $E(\mathbf w)=\lambda$，
则存在 $\mathbf w_0\in W$，
使 $T^\lambda\mathbf w_0=\mathbf w$。
那么 $W$ 有限生成，
且有一组元素能量全为零的标准基。
\end{lemma}

\begin{proof}
令 $W_{00}$ 为 $\overline V$ 中
可作为 $W$ 某个元素首系数的元素之集。
$W_{00}\cup\{0\}$ 显然是
$\overline V$ 的有限维 $R$-线性子空间，
记为 $W_0$。
取 $W_0$ 的一组基 $\{w_i\}$，
对每个 $w_i$ 选取
$\mathbf w_i\in W$，
使 $\sigma(\mathbf w_i)=w_i$。
由引理的假设可设
$E(\mathbf w_i)=0$。
我们证明 $\{\mathbf w_i\}$ 生成 $W$；
由此即得引理。

对 $V$ 施以滤过自同构后，
不妨设
\begin{equation}
\mathbf w_i=w_i\in\overline V.
\tag{6.3.12}\label{eq:6-3-12}
\end{equation}
现在取 $\mathbf v\in W$，写为
\[
\mathbf v=\sum_iT^{\lambda_i}v_i.
\]
依定义 $v_1\in W_0$；
由 (6.3.12) 可知 $v_1\in W$，
因而
$\sum_{i=2}^{\infty}T^{\lambda_i}v_i\in W$。
于是 $v_2\in W_0$。
重复同一论证，
可知每个 $v_i\in W_0$。
由于每个 $v_i$ 都是 $w_i$ 的
$R$-线性组合，
$\mathbf v$ 属于 $\mathbf w_i$
生成的 $\Lambda_{0,\mathrm{nov}}^{(0)}$-子模。
\end{proof}

\pagebreak
\begin{corollary}{6.3.13}\label{cor:6-3-13}
设 $(C,\delta)$ 是有限 d.g.c.f.z.。
则对每个 $k$，
$\ker\delta_k=\ker\delta\cap C^k$
都是有限生成的。
而且，可为它们取全部由能量零元素
组成的标准基。
\end{corollary}

\begin{proof}
容易看出 $\ker\delta$
满足引理~\ref{lem:6-3-11} 的假设。
\end{proof}

推论~\ref{cor:6-3-13}
与引理~\ref{lem:6-3-2}
立即给出有限 d.g.c.f.z.
的上同调结构定理。

\begin{proposition}{6.3.14}\label{prop:6-3-14}
设 $(C,\delta)$ 是有限 d.g.c.f.z.。
则对每个整数 $p$，
存在有限多个常数
$\lambda_{p,i}>0$（$i=1,\ldots,k_p$）
及整数 $\ell_p$，
其中 $\ell_p\le\operatorname{rank}C^p-k_p$，
使
\begin{equation}
H^p(C,\delta)
\cong
\left(
\bigoplus_{i=1}^{k_p}
\Lambda_{0,\mathrm{nov}}^{(0)}
/F^{\lambda_{p,i}}\Lambda_{0,\mathrm{nov}}^{(0)}
\right)
\oplus
\bigl(\Lambda_{0,\mathrm{nov}}^{(0)}\bigr)^{\oplus\ell_p}.
\tag{6.3.15}\label{eq:6-3-15}
\end{equation}
若按 $\lambda_{p,i+1}\ge\lambda_{p,i}$ 排序，
这些 $\lambda_{p,i}$ 唯一确定。
\end{proposition}

\begin{proof}
先注意，$C^{p-1}$ 有限，
故 $\operatorname{Im}\delta_{p-1}$
有限生成并有标准基，
可写为
\[
\{T^{\mu_{p,1}}u_1,
T^{\mu_{p,2}}u_2,\ldots,
T^{\mu_{p,m_p}}u_{m_p}\},
\]
其中 $u_i$ 是 $\ker\delta_p$
中的能量零元素，
其约化属于 $\overline C^p
=C^p/\Lambda_{0,\mathrm{nov}}^{+,(0)}C^p$。
事实上，像中的标准基元可提取其
$T^{\mu_{p,i}}$ 因子；
由 $\delta_p(T^{\mu_{p,i}}u_i)=0$
及 $C^{p+1}$ 无挠，
即得 $\delta_pu_i=0$。
\translatornote{原书实体页 380
在此先写 $\operatorname{Im}\delta_p$
及“因 $C^p$ 有限”，
随后却使用
$\operatorname{Im}\delta_{p-1}\subset\ker\delta_p$；
这里按型别改用前一阶微分的像，
并把 $u_i$ 明确为能量零的提升，
见疑误 A-0365。}
由
$\operatorname{Im}\delta_{p-1}
\subset\ker\delta_p$，
以及 $\ker\delta_p$
有能量零标准基，
可将 $\{u_1,\ldots,u_{m_p}\}$
扩充为 $\ker\delta_p$ 的标准基
\[
\{u_1,\ldots,u_{m_p},
u_{m_p+1},\ldots,u_{m_p+\ell_p}\}.
\]
由此立即得到 (6.3.15)：
对所有 $\mu_{p,i}>0$，
取 $\lambda_{p,i}=\mu_{p,i}$，
而 $\mu_{p,i}=0$ 对应的零模不列入直和。\translatornote{原书命题要求
$\lambda_{p,i}>0$，
但证明未交代像的标准基若有零能量元素
该如何处理；此类元素对应的商模为零，
故在最终直和中略去，
见疑误 A-0366。}

为证唯一性，只需注意
\[
\dim_R
\frac{T^\lambda H^p(C,\delta)}
{T^\lambda\Lambda_{0,\mathrm{nov}}^{+,(0)}
H^p(C,\delta)}
=\ell_p+\#\{i\mid\lambda_{p,i}>\lambda\}.
\]
\end{proof}

作者尚不知道，
若把域 $R$ 换成任意诺特环，
命题~\ref{prop:6-3-9}
（因而谱序列的收敛性）是否仍成立。
